\documentclass[../../../include/open-logic-section]{subfiles}

\begin{document}

\olfileid{sth}{replacement}{strength}
\olsection{قدرت جایگزینی}

با مشاهده‌ای ساده دربارهٔ قدرت جایگزینی آغاز می‌کنیم: تا هنگامی که از
$\Z$ فراتر نرویم، نمی‌توانیم وجود هیچ عدد ترتیبی فون نویمانِ بزرگ‌تر یا
مساویِ~$\omega+\omega$ را اثبات کنیم.

طرح دلیل چنین است. در~$\ZF$ کار کنید و مجموعهٔ~$V_{\omega+\omega}$ را
در نظر بگیرید. این مجموعه دامنهٔ \emph{مدلی} برای~$\Z$ است. برای دیدن
این امر، نمادگذاری‌ای برای \emph{نسبی‌سازیِ} یک فرمول معرفی می‌کنیم:
\begin{defn}\ollabel{formularelativization}
	برای هر مجموعهٔ~$M$ و هر فرمول~$\phi$، بگذارید~$\phi^M$ فرمولی باشد
	که با محدودکردن همهٔ سورهای~$\phi$ به~$M$ حاصل می‌شود. یعنی
	«$\lexists[x]$» را با «$(\lexists[x \in M])$» و
	«$\lforall[x]$» را با «$(\lforall[x \in M])$» جایگزین کنید.
\end{defn}\noindent
می‌توان نشان داد برای هر اصل موضوع~$\phi$ از~$\Z$،
$\ZF\vdash\phi^{V_{\omega+\omega}}$. اما بنا بر
\olref[spine][rank]{ordsetrankalpha}، $\omega+\omega$ در
$V_{\omega+\omega}$ \emph{نیست}. پس~$\Z$ با نبودِ~$\omega+\omega$
سازگار است.

از همین رو در~\olref[ordinals][replacement]{sec} گفتیم
\olref[ordinals][ordtype]{thmOrdinalRepresentation} را نمی‌توان بدون
جایگزینی اثبات کرد. زیرا درون~$\Z$ به‌آسانی می‌توان خوش‌ترتیبیِ صریحی
تعریف کرد که به‌طور شهودی \emph{باید} نوع ترتیب~$\omega+\omega$ داشته
باشد. در واقع، در~\olref[ordinals][idea]{sec} نمونه‌ای غیرصوری از آن
دادیم: ترتیب زیر روی اعداد طبیعی:
\begin{align*}
	n \lessdot m \text{ اگر و تنها اگر }&\text{یا }n < m\text{ و }m-n\text{ زوج باشد،}\\
	& \text{یا $n$ زوج و $m$ فرد باشد.}
\end{align*}
اما اگر~$\omega+\omega$ وجود نداشته باشد، این خوش‌ترتیب با هیچ عدد
ترتیبی‌ای هم‌ریخت نیست. پس~$\Z$،
\olref[ordinals][ordtype]{thmOrdinalRepresentation} را \emph{اثبات
نمی‌کند}.

اکنون مطلب را وارونه بنگریم: جایگزینی به ما اجازه می‌دهد وجود
$\omega+\omega$ و در نتیجه وجود~$V_{\omega+\omega}$ را اثبات کنیم.
اما موضوع به همین‌جا ختم نمی‌شود. برای \emph{هر} خوش‌ترتیبی که بتوانیم
تعریف کنیم، \olref[ordinals][ordtype]{thmOrdinalRepresentation} می‌گوید
یک~$\alpha$ وجود دارد که با آن خوش‌ترتیب هم‌ریخت است، و بنابراین
$V_\alpha$ وجود دارد. پس جایگزینی، به‌شیوه‌ای سرراست، تضمین می‌کند که
سلسله‌مراتب مجموعه‌ها باید \emph{بسیار بلند} باشد.

در چند بخش آینده و سپس بار دیگر در
\olref[card-arithmetic][fix]{sec} درک بهتری از اینکه جایگزینی
سلسله‌مراتب را تا چه اندازه بلند می‌کند به دست خواهیم آورد. فعلاً نکتهٔ
ساده این است که جایگزینی واقعاً به توجیه \emph{نیاز دارد}!

\end{document}
